% Standalone section draft for insertion into the ICLR manuscript.
% Requires booktabs and natbib. Citation keys are in paper/references.bib.
% This file does not modify or automatically enter main.tex.
% Sources: current manuscript; docs/worldmem_final_evaluation_handoff.md;
% VMem transfer audit dated 2026-09-06 supplied by the user;
% matched-cohort VBench diagnosis dated 2026-09-14; fresh GeoCov VBench run.
% MemCam 60s GeoCov LPIPS/FVD below are previously recorded values.
% Final source-file/config provenance for that pair still requires checking.
% No 30-video stored VBench aggregates are used here.
% Add the VBench++ reference before submission; Long citation is TBD below.

\providecommand{\Geo}{\textsc{GeoCov}}
\providecommand{\RI}{\textsc{RI}}
\providecommand{\TBD}{\textbf{TBD}}

\section{Experiments}
\label{sec:experimental-results}

Our experiments ask whether a small, maintained archive can replace growing
history without sacrificing the evidence needed for camera-controlled video
generation. We evaluate three aspects of this design: video quality at a fixed
memory budget, transfer across memory representations and retrieval mechanisms,
and the tradeoff between retained history and the evidence actually retrieved.
MemCam and WorldMem provide the main quantitative comparisons. A VMem-based
transfer experiment tests whether archive curation also complements
surfel-indexed retrieval.

\subsection{Experimental Setup}
\label{sec:exp-protocol}

\paragraph{MemCam.}
We use the public MemCam implementation with the Wan2.1 1.3B backbone
\citep{memcam,wan}. The primary suite contains fifteen matched 60-second
trajectories from the Context-as-Memory dataset \citep{contextasmemory}.
All policies share the input images, prompts, camera trajectories, generation
seeds, model checkpoint, and 50-step denoising schedule. Generation proceeds
in sections of 76 target frames. We evaluate $B\in\{16,32,64,128\}$, with
$B=32$ fixed for the main policy comparison. A separate fifteen-trajectory
180-second experiment tests longer rollouts; its scores are not pooled with
the 60-second suite. The controller changes the persistent archive update,
while leaving the generator and FOV-based retriever unchanged.

\paragraph{WorldMem.}
We evaluate fifteen matched 60-second videos using WorldMem's latent-memory
interface \citep{worldmem}. All policies are evaluated on the same first
fifteen batch identities, including when an output directory contains additional
videos. We compare complete retention and the same five bounded policy families
at budgets of 16, 32, 64, and 128 items. The geometric controller is adapted to
WorldMem's memory representation, rather than applying the decoded-RGB MemCam
implementation unchanged. The checkpoint, dataset split, frame rate,
resolution, denoising configuration, exact latent-item definition, and visual
descriptor used by this adaptation are \TBD. Comparisons are made within
WorldMem, not between absolute metric values from different generators.

\paragraph{VMem transfer.}
The transfer protocol uses a VMem-based generator with surfel-indexed view
retrieval \citep{vmem}. It pairs unbounded memory with adapted \Geo-32 on
five input images, each with three camera programs: two angular sweeps
($-45^\circ$ to $45^\circ$ and $-90^\circ$ to $90^\circ$) and an out-and-back
translation. This gives fifteen paired cases, not fifteen independent scenes.
Each case uses 195 actions and produces 781 frames at 13 fps, approximately
60 seconds, at $576\times576$ resolution. The two arms share the seed,
checkpoint, four reference and four target slots, 50 denoising steps, and
400 reconstruction iterations. Confirmed completed pairs: \TBD/15.

VMem supplies CLIP image embeddings, with a latent fallback, in place of
MemCam's DINO descriptors; the geometric-policy coefficients remain unchanged.
Both arms use the same audited VMem fork, including its shared retrieval and
fallback modifications. In the budgeted arm, eviction also removes references
from the surfel index, deletes surfels with no remaining references, and changes
the frames available for later reconstruction. This is therefore an end-to-end
memory-controller transfer, not a retrieval-only intervention. Here $B$ bounds
eligible views; it does not bound all output buffers, all historical state,
or the number of surfels.

\paragraph{Comparison policies.}
Complete retention keeps all historical items. FIFO prioritizes recency;
K-center prioritizes descriptor coverage; Marginal Coverage Eviction (MCE)
uses a marginal coverage objective; and rarity--irreplaceability (\RI)
prioritizes rare appearance groups without close substitutes. \Geo instead
scores redundancy jointly in camera-pose and appearance space. The main
MemCam and WorldMem comparison fixes the bounded policies at $B=32$; the budget
sweep separates the effect of policy choice from capacity. The VMem transfer
uses only the paired unbounded and \Geo-32 arms.
\Geo protects the initial conditioning frame and current section endpoint;
the baseline-by-baseline protection configuration and budget accounting for
MemCam and WorldMem are \TBD. In the audited VMem adaptation, both protected
views count toward $B$.

\paragraph{Metrics and matching.}
LPIPS \citep{lpips} measures fidelity against trajectory-indexed ground truth.
The LPIPS network and spatial preprocessing are \TBD. FVD \citep{fvd} uses
16-frame clips, four clips per video, frame stride four, and 224-pixel inputs,
with the same I3D detector within each comparison. The exact clip-start rule
and detector/configuration identifiers are \TBD. We report FVD as a
distribution-level statistic, not an average of per-video FVD scores.

We additionally report six reference-free VBench dimensions \citep{vbench}:
subject consistency, background consistency, motion smoothness, dynamic degree,
aesthetic quality, and imaging quality. VBench-Long is evaluated separately
using its custom-video protocol and final dimension scores (benchmark citation:
\TBD). Each policy is evaluated on the same original fifteen videos, without
the extra seeds present in some generation directories. Per-video scores are
averaged with equal video weight. We report the dimensions separately rather
than constructing a new overall quality score. These appearance-based scores
complement, rather than replace, trajectory-conditioned fidelity evaluation.

CUT3R-based camera evaluation \citep{cut3r} is conditional on passing a
ground-truth camera sanity check. Validated rotation and translation results
are currently \TBD; provisional outputs from the unresolved calibration are
not treated as camera-accuracy evidence. The VMem demo protocol has no dense
ground-truth novel-view video, so it does not support the same GT-based
LPIPS/FVD comparison as the main suites.

\paragraph{Resource evaluation.}
We distinguish the number of retained items from archive bytes, total host
memory, peak GPU memory, and retrieval latency. Timing measurements must use
matched hardware and generation settings, CUDA synchronization where relevant,
and a stated warm-up policy. The device, timing protocol, and measured results
are \TBD. End-to-end timing includes descriptor extraction and archive updates;
retrieval-only timing is reported separately.

\subsection{Video Quality at a Fixed Budget}
\label{sec:exp-quality}

\begin{table}[t]
\centering
\small
\caption{Primary 60-second evaluation on fifteen matched videos per system.
All bounded rows use $B=32$. Lower is better. \TBD\ denotes an unfilled
result, not a zero or a failed run. Absolute scores are compared only within
each system.}
\label{tab:exp-quality-60}
\begin{tabular}{lrrrr}
\toprule
& \multicolumn{2}{c}{MemCam} & \multicolumn{2}{c}{WorldMem} \\
\cmidrule(lr){2-3}\cmidrule(lr){4-5}
Policy & LPIPS & FVD & LPIPS & FVD \\
\midrule
Complete retention & \TBD & \TBD & 0.652 & 3077.6 \\
FIFO & \TBD & \TBD & 0.689 & 3554.9 \\
\RI / Latent-\RI & \TBD & \TBD & 0.546 & 1160.4 \\
K-center & \TBD & \TBD & 0.559 & \TBD \\
MCE & \TBD & \TBD & 0.575 & \TBD \\
\Geo & 0.5850 & 690.5 & 0.534 & 1116.9 \\
\bottomrule
\end{tabular}
\end{table}

Table~\ref{tab:exp-quality-60} compares retention policies under a common
budget. In WorldMem, \Geo-32 lowers LPIPS from 0.652 to 0.534 and FVD from
3077.6 to 1116.9 relative to complete retention. FIFO-32 worsens both metrics,
while Latent-\RI improves both. Among the reported B32 policies, \Geo has
the lowest LPIPS, and the lowest FVD among policies with available FVD values.
The result supports structured archive composition rather than recency alone.
The corresponding complete MemCam 60-second comparison is \TBD.

\paragraph{Extension to 180 seconds.}
The longer MemCam experiment tests whether the bounded archive remains useful
as the rollout grows. Table~\ref{tab:exp-quality-180} shows that \Geo-32
retains 32 rather than 5,397 frames while lowering FVD from 734.2 to 476.6
(35.1\%) and LPIPS from 0.5980 to 0.5876. FIFO improves FVD but worsens
trajectory-indexed LPIPS. Thus, a lower distributional video score alone does
not establish better reconstruction of the commanded scene.

\begin{table}[t]
\centering
\small
\caption{MemCam duration extension: fifteen matched 180-second videos.
Lower LPIPS and FVD are better.}
\label{tab:exp-quality-180}
\begin{tabular}{lrrr}
\toprule
Policy & Retained frames & LPIPS & FVD \\
\midrule
Complete retention & 5,397 & 0.5980 & 734.2 \\
FIFO-32 & 32 & 0.6514 & 677.3 \\
\RI-32 & 32 & 0.5939 & 550.4 \\
\Geo-32 & 32 & \textbf{0.5876} & \textbf{476.6} \\
\bottomrule
\end{tabular}
\end{table}

\subsection{Perceptual Quality and Temporal Consistency}
\label{sec:exp-vbench}

\begin{table}[t]
\centering
\small
\caption{MemCam, fifteen matched 60-second videos. Scores are multiplied by
100; higher is better according to each metric. Standard VBench and
VBench-Long are separate evaluators. RI's incomplete Long cohort is not
reported as a fifteen-video mean.}
\label{tab:exp-vbench}
\begin{tabular}{lrrrrrr}
\toprule
Policy & Subject & Background & Motion & Dynamic & Aesthetic & Imaging \\
\midrule
\multicolumn{7}{l}{\emph{Standard VBench}} \\
Complete retention & 80.53 & 90.17 & 99.27 & 100.00 & 45.42 & 49.25 \\
FIFO-32 & 77.01 & 89.51 & 99.28 & 100.00 & 45.35 & 50.78 \\
\RI-32 & 81.28 & 89.80 & 99.19 & 100.00 & 45.99 & 50.16 \\
\Geo-32 & 81.78 & 90.47 & 99.21 & 100.00 & 46.14 & 50.54 \\
\midrule
\multicolumn{7}{l}{\emph{VBench-Long}} \\
Complete retention & 95.26 & 95.52 & 99.28 & 93.12 & 43.45 & 47.63 \\
FIFO-32 & 95.47 & 95.58 & 99.30 & 91.18 & 43.49 & 48.95 \\
\RI-32 & \TBD & \TBD & \TBD & \TBD & \TBD & \TBD \\
\Geo-32 & 94.70 & 95.36 & 99.23 & 97.42 & 44.13 & 48.81 \\
\bottomrule
\end{tabular}
\end{table}

On standard VBench, \Geo-32 has the highest subject consistency, background
consistency, and aesthetic quality among the four displayed policies
(Table~\ref{tab:exp-vbench}). Its imaging quality improves over complete
retention but is slightly below FIFO; motion smoothness is also slightly
lower than FIFO. Dynamic degree is saturated for all four policies.
The fresh GeoCov-only evaluation reproduces the previously matched standard
scores closely; the table uses the fresh GeoCov values and the matched
historical values for the other policies.

VBench-Long exposes a different tradeoff. \Geo has higher dynamic degree and
aesthetic quality than complete retention and FIFO, but lower subject and
background consistency. Its imaging quality exceeds complete retention but
not FIFO. These results do not support uniform superiority across temporal
consistency metrics. In the paired diagnostic, the subject-consistency
disadvantage relative to FIFO is concentrated in the within-clip component;
the positive raw cross-clip background difference is attenuated by the
benchmark's score mapping before fusion. This decomposition explains how the
reported ranking is formed, but does not replace the final benchmark scores
or establish a causal explanation for the videos' behavior.

The corresponding WorldMem VBench and VBench-Long results are \TBD.
Complete matched RI VBench-Long coverage remains \TBD/15; the supplied
evaluation currently contains seven completed videos.

\subsection{Transfer to Surfel-Indexed Memory}
\label{sec:exp-vmem}

VMem tests a complementary setting: the reader already uses spatial indexing
and geometric reference selection. The comparison asks whether maintaining
the eligible archive adds value on top of that reader, rather than replacing
it. Table~\ref{tab:exp-vmem} reserves the matched reference-free evaluation.
The measured quality effect, per-case variation, and completed-pair count are
\TBD. No positive transfer claim is inferred from the availability of the
implementation alone.

\begin{table}[t]
\centering
\small
\caption{VMem transfer protocol: fifteen planned paired 60-second cases
across five input images. Results require completed-pair and configuration
verification. Scores are on the same $\times100$ scale as
Table~\ref{tab:exp-vbench}.}
\label{tab:exp-vmem}
\begin{tabular}{lrrrr}
\toprule
& \multicolumn{2}{c}{VBench} & \multicolumn{2}{c}{VBench-Long} \\
\cmidrule(lr){2-3}\cmidrule(lr){4-5}
Dimension & Unbounded & \Geo-32 & Unbounded & \Geo-32 \\
\midrule
Subject consistency & \TBD & \TBD & \TBD & \TBD \\
Background consistency & \TBD & \TBD & \TBD & \TBD \\
Motion smoothness & \TBD & \TBD & \TBD & \TBD \\
Dynamic degree & \TBD & \TBD & \TBD & \TBD \\
Aesthetic quality & \TBD & \TBD & \TBD & \TBD \\
Imaging quality & \TBD & \TBD & \TBD & \TBD \\
\bottomrule
\end{tabular}
\end{table}

\subsection{Budget Sensitivity and Resource Scaling}
\label{sec:exp-budget}

The budget sweep tests whether the B32 result depends on a single capacity
choice. In WorldMem, \Geo's 60-second LPIPS is 0.525, 0.534, 0.545, and
0.577 at B16, B32, B64, and B128, respectively, compared with 0.652 for
complete retention. In the MemCam 180-second extension, the corresponding
LPIPS values are 0.5876, 0.5876, 0.5865, and 0.5913, and FVD values are
446.1, 476.6, 493.9, and 515.5. The primary MemCam 60-second budget curve
is \TBD. These within-policy observations show that additional retained
history does not necessarily improve final quality; they do not establish
a universal optimal budget or an intrinsic penalty for candidate count.

For MemCam's exact retriever, an archive receiving $s$ items per chunk stores
$O(Ts)$ items after $T$ chunks and requires $O(T^2s)$ cumulative candidate
comparisons for a fixed number of queries per chunk. A bounded archive uses
$O(B)$ persistent items and $O(TB)$ cumulative comparisons. The graph update
cost is $O((B+s)^2)$ per chunk in the direct implementation. These are
archive and candidate-scoring bounds, not constant-total-memory or measured
end-to-end speed claims. Approximate indexing may accelerate a growing
archive without bounding its storage.

Measured archive bytes, retrieval latency, update overhead, host memory,
and peak GPU memory are \TBD\ for the matched duration sweep. WorldMem's
representation-specific measurements are \TBD. For VMem, report eligible-view
count and surfel count separately: the current adapter bounds the former,
while the latter and additional retained process state may still grow.

\subsection{What Makes the Retained Archive Useful?}
\label{sec:exp-mechanism}

The memory diagnostics separate retaining an informative observation from
actually retrieving it. On thirteen common-source trajectories at B32,
complete retention has zero retention gap but a selection gap of 0.2188.
\Geo incurs a retention gap of 0.0426 while reducing the selection gap to
0.1470. \RI retains slightly better hindsight evidence (0.0399 retention
gap), but leaves more of its quality unused by the retriever (0.1609
selection gap). The diagnostics therefore measure different consequences
of archive design rather than treating retention capacity as retrieval quality.

A separate common-source analysis fixes candidate pixels to one rollout and
compares selected images against ground truth at their own historical indices.
Across fifteen trajectories, \Geo's selected-frame PSNR exceeds complete
retention by 4.629 dB and SSIM by 0.1512. This measures the fidelity of selected
historical content, not whether its view exactly matches the current query.
Both archives contain the initial conditioning image. Excluding disagreements
where \Geo selects that image leaves a 1.060 dB PSNR advantage and 0.0449
SSIM advantage. Restricting further to selections with FOV overlap at least
0.80 in both policies gives differences of 0.103 dB and 0.0098, with both
confidence intervals crossing zero. Thus, the large aggregate effect is not
evidence that \Geo generally detects cleaner generated images at matched views.

Finally, randomly expanding a fixed recent-history B32 pool to full history
changes the retrieval winner on 72.2\% of queries, but does not significantly
reduce selected-frame fidelity. The full-history minus B32 differences are
0.228 dB PSNR (95\% CI $[-0.109,0.664]$) and 0.0011 SSIM
(95\% CI $[-0.0117,0.0188]$). This intervention does not support
candidate-count-only degradation in the tested setting. Together, the
controls motivate archive composition as a design variable without claiming
that archive growth alone causes a snowballing loss of quality.

\paragraph{Scope of the evidence.}
Policy development used the MemCam suite; it is not an untouched test set.
LPIPS and FVD results above are point estimates, with matched uncertainty
\TBD. Existing paired VBench diagnostic intervals bootstrap trajectories
and are exploratory, without multiple-comparison adjustment. For VMem,
uncertainty should account for the five shared scene inputs rather than
treating all fifteen camera cases as independent scenes. The supported
conclusion is that structured bounded retention can improve the reported
generation and fidelity metrics, while temporal-consistency tradeoffs,
validated camera accuracy, and measured system-level savings remain
separate empirical questions.
